\scanpage{104}
\mspage{19}
% archive sheet number: 103
% source lines 1--6, top additions and cancellations
\uuline{Démonstration} \deleted{Pour (4.2.1), la question étant locale,} \deleted{la \unsure{pureté} \unsure{par|pour}} \inserted{Prouvons d'abord (4.2.1) et le fait que} \deleted{\unsure{le}} $\underline H^{2d}_Y(\Lambda_X)$ est loc. isom. à $\Lambda_Y$. \deleted{\unsure{raisonnons sur $d$,}} Le cas $d=0$ est trivial, le cas $d=1$ résulte de la théorie de Kummer et du \unsure{Lemme|lemme} de Abhyankar absolu, \deleted{\unsure{déduisant}} \inserted{qui fournit $\underline H^2_Y(\mu_{nX})\simeq\Lambda_Y$} (\deleted{SGA 4} cf SGA 4 $\overline{\underline{\mathrm{XIX}}}$ 1.3 ou SGA 1 $\overline{\underline{\mathrm{XII}}}$ \ldots).

% source lines 7--11
Pour $d>1$, on procède par récurrence sur $d$. L'assertion étant locale, on \deleted{\unsure{se ramène}} est ramené au cas où il existe un \uline{diviseur régulier} $Z$ dans $X$ tel que \deleted{\unsure{$Z$}} $Y\subset Z$. \deleted{On utilise} \inserted{Considérons} La suite spectrale de transitivité

\[
\underline H^*_Y(\Lambda_X)\Longleftarrow E^{pq}_2=\underline H^p_Y(\underline H^q_Z(\Lambda_X)),
\]

% source lines 13--17
\deleted{et l'hyp. de réc.} alors \deleted{\unsure{$E^p$}} $\underline H^q_Z(\Lambda_X)$ est nul si $q\ne 2$, et loc. isom. à $\Lambda_Y$ si $\unsure{q|p}=1$, \unsure{donc|dans} l'hyp. de réc. nous prouve que $E^{pq}_2=0$ sauf \deleted{\unsure{p.}} si \deleted{$p=d-1$,} $p=2(d-1)$, $q=2$, et que pour ces valeurs de $p,q$, $E^{pq}_2$ est loc. isom. à $\Lambda_Z$, d'où la conclusion.

% source lines 18--22
Pour la définition d'un isom. canonique (4.2.1) ou (4.2.1 bis), le cas $d=1$ est traité dans (SGA 4 XIX 1.3), \deleted{le cas $d$} Pour le cas $d$ quelconque, il suffit d'utiliser la construction de (SGA 4 $\overline{\underline{\mathrm{XIX}}}$ 3.4) ou de (SGA 4 $\overline{\underline{\mathrm{XVIII}}}$\quad).
